\documentclass[10pt,a4paper]{amsart}
\usepackage{amsmath,amssymb,amsthm,graphicx}
\usepackage{verbatim}
\newenvironment{reference}{\comment}{\endcomment}
%\newenvironment{reference}{}{}
\newenvironment{slogan}{\comment}{\endcomment}
\newenvironment{history}{\comment}{\endcomment}

% For commutative diagrams we use Xy-pic
\usepackage[all]{xy}

% We use 2cell for 2-commutative diagrams.
\xyoption{2cell}
% Theorem environments.
%
\theoremstyle{plain}
\newtheorem{theorem}[subsection]{Theorem}
\newtheorem{proposition}[subsection]{Proposition}
\newtheorem{lemma}[subsection]{Lemma}

\theoremstyle{definition}
\newtheorem{definition}[subsection]{Definition}
\newtheorem{example}[subsection]{Example}
\newtheorem{exercise}[subsection]{Exercise}
\newtheorem{situation}[subsection]{Situation}

\theoremstyle{remark}
\newtheorem{remark}[subsection]{Remark}
\newtheorem{remarks}[subsection]{Remarks}

\numberwithin{equation}{subsection}

% Macros
%
\def\lim{\mathop{\mathrm{lim}}\nolimits}
\def\colim{\mathop{\mathrm{colim}}\nolimits}
\def\Spec{\mathop{\mathrm{Spec}}}
\def\Hom{\mathop{\mathrm{Hom}}\nolimits}
\def\Ext{\mathop{\mathrm{Ext}}\nolimits}
\def\SheafHom{\mathop{\mathcal{H}\!\mathit{om}}\nolimits}
\def\SheafExt{\mathop{\mathcal{E}\!\mathit{xt}}\nolimits}
\def\Sch{\mathit{Sch}}
\def\Mor{\mathop{\mathrm{Mor}}\nolimits}
\def\Ob{\mathop{\mathrm{Ob}}\nolimits}
\def\Sh{\mathop{\mathit{Sh}}\nolimits}
\def\NL{\mathop{N\!L}\nolimits}
\def\CH{\mathop{\mathrm{CH}}\nolimits}
\def\proetale{{pro\text{-}\acute{e}tale}}
\def\etale{{\acute{e}tale}}
\def\QCoh{\mathit{QCoh}}
\def\Ker{\mathop{\mathrm{Ker}}}
\def\Im{\mathop{\mathrm{Im}}}
\def\Coker{\mathop{\mathrm{Coker}}}
\def\Coim{\mathop{\mathrm{Coim}}}

% Boxtimes
%
\DeclareMathSymbol{\boxtimes}{\mathbin}{AMSa}{"02}

%
% Macros for moduli stacks/spaces
%
\def\QCohstack{\mathcal{QC}\!\mathit{oh}}
\def\Cohstack{\mathcal{C}\!\mathit{oh}}
\def\Spacesstack{\mathcal{S}\!\mathit{paces}}
\def\Quotfunctor{\mathrm{Quot}}
\def\Hilbfunctor{\mathrm{Hilb}}
\def\Curvesstack{\mathcal{C}\!\mathit{urves}}
\def\Polarizedstack{\mathcal{P}\!\mathit{olarized}}
\def\Complexesstack{\mathcal{C}\!\mathit{omplexes}}
% \Pic is the operator that assigns to X its picard group, usage \Pic(X)
% \Picardstack_{X/B} denotes the Picard stack of X over B
% \Picardfunctor_{X/B} denotes the Picard functor of X over B
\def\Pic{\mathop{\mathrm{Pic}}\nolimits}
\def\Picardstack{\mathcal{P}\!\mathit{ic}}
\def\Picardfunctor{\mathrm{Pic}}
\def\Deformationcategory{\mathcal{D}\!\mathit{ef}}
\usepackage{cleveref}
\crefname{originalremark}{original remark}{original remarks}
\Crefname{originalremark}{Original remark}{Original remarks}
\makeatletter
\newlabel{spaces-properties-remark-explain-equivalence}{{03H7}{source}{}{sourcepointer-1}{}}
\newlabel{spaces-properties-remark-explain-equivalence@cref}{{[originalremark][0][]03H7}{[1][0][]source}{}{}{}}
\makeatother
\title{An exact sequence computing the pushforward along a morphism of algebraic spaces}
\author{The Stacks Project Authors}
\date{stacks-project:spaces-morphisms.tex:local-snapshot-20261012; GNU FDL 1.2 or later}
\begin{document}
\maketitle
\setcounter{subsection}{1}
\noindent\textit{Editorial note.} All delivered mathematical text is the author's own, copied from the pinned Stacks Project chapter file spaces-morphisms.tex; no statement or proof step was rewritten, reordered or removed.\par\medskip
\noindent\textit{Editorial note.} The displayed exact sequence in the statement carries the sentence identifying the second arrow as the difference t\textasciicircum{}* - s\textasciicircum{}*; the proof shows the difference arises from the sheaf condition for the covering V x\_Y U -> V x\_Y X.\par\medskip
\noindent\textit{Editorial note.} The single reference of the counted spans is recorded as a non-load-bearing navigation pointer; the argument that consumes it stays inside the counted proof span.\par\medskip
\noindent\textit{Editorial note.} Preamble: the three native excerpt spans of sources/stacks-spaces-morphisms/preamble.tex (lines 8-13, 14-18 and 154-217). The chapter file itself carries no preamble text.\par\medskip
\noindent\textit{Allowed standard prerequisites (editorial).}\begin{itemize}
\item Sheaves of modules on the big etale site and the equivalence with modules on X\_etale. Properties of Spaces, Remark spaces-properties-remark-explain-equivalence, used at proof lines 2069-2073. Justifies evaluating F on objects of X\_\{spaces,etale\}.
\item The sheaf condition of an O\_X-module on the big etale site. used at proof lines 2074-2086. Produces the short exact sequence of sections over V.
\item Identification of the fibre product (V x\_Y U) x\_X (V x\_Y U) with V x\_Y R. used at proof lines 2078-2079. Turns the sheaf condition into the claimed sequence with the difference of t\textasciicircum{}* and s\textasciicircum{}*.
\end{itemize}\medskip
\ifdefined\hypertarget\hypertarget{sourcepointer-1}{}\fi
\noindent\textit{Source reference.} Original remark 03H7: focused-seeds/AGP/sources/stacks-snapshot-20261011/tags.txt, source lines 4556--4556. Navigation pointer into the Stacks Project; the standard fact it supplies is applied inside the counted proof span and the argument resting on it stays there.\par\medskip
\begin{lemma}
\label{lemma-compute-pushforward}
Let $S$ be a scheme.
Let $f : X \to Y$ be a morphism of algebraic spaces over $S$.
Let $U \to X$ be a surjective \'etale morphism from a scheme to $X$.
Set $R = U \times_X U$ and denote $t, s : R \to U$ the projection
morphisms as usual. Denote $a : U \to Y$ and $b : R \to Y$ the induced
morphisms. For any object $\mathcal{F}$ of $\textit{Mod}(\mathcal{O}_X)$
there exists an exact sequence
$$
0 \to f_*\mathcal{F} \to a_*(\mathcal{F}|_U) \to b_*(\mathcal{F}|_R)
$$
where the second arrow is the difference $t^* - s^*$.
\end{lemma}

\begin{proof}
We denote $\mathcal{F}$ also its extension to a sheaf of modules on
$X_{spaces, \etale}$, see
Properties of Spaces,
Remark \ref{spaces-properties-remark-explain-equivalence}.
Let $V \to Y$ be an object of $Y_\etale$. Then $V \times_Y X$ is an
object of $X_{spaces, \etale}$, and by definition
$f_*\mathcal{F}(V) = \mathcal{F}(V \times_Y X)$. Since $U \to X$ is
surjective \'etale, we see that $\{V \times_Y U \to V \times_Y X\}$
is a covering. Also, we have
$(V \times_Y U) \times_X (V \times_Y U) = V \times_Y R$.
Hence, by the sheaf condition of $\mathcal{F}$ on
$X_{spaces, \etale}$ we have a short exact sequence
$$
0 \to \mathcal{F}(V \times_Y X)
\to \mathcal{F}(V \times_Y U) \to \mathcal{F}(V \times_Y R)
$$
where the second arrow is the difference of restricting via $t$ or $s$.
This exact sequence is functorial in $V$ and hence we obtain the lemma.
\end{proof}

\end{document}
